%=====================================================================
\chapter{Taking Payment for Orders}\label{ch:gateway}
%=====================================================================
\locator{4}{Part IV --- Money and Identity}

\begin{outcomes}
By the end of this chapter you will be able to:
\begin{tight}
  \item \textbf{Describe} the handshake between a shop and a gateway, and
        \textbf{name} the three outcomes of a payment request.
  \item \textbf{Measure} a gateway's response-time distribution, and
        \textbf{choose} a timeout from it rather than from its median.
  \item \textbf{Explain} why a timed-out payment has an unknown outcome, and
        \textbf{design} what the shop does next.
  \item \textbf{Implement} idempotency, and \textbf{demonstrate} what happens
        without it.
  \item \textbf{State} what PCI DSS requires of a shop that never touches a
        card number, and \textbf{identify} what would change that.
\end{tight}
\end{outcomes}

\begin{prereq}
\chref{payments}: the parties, and the difference between authorisation,
capture and settlement. \chref{checkout}: the order state machine, status 0,
and why the order identifier exists before the payment.
\end{prereq}

\begin{keyterms}
gateway integration $\bullet$ hosted page $\bullet$ direct post $\bullet$
server-to-server $\bullet$ callback $\bullet$ webhook $\bullet$ timeout
$\bullet$ latency distribution $\bullet$ percentile $\bullet$ long tail
$\bullet$ idempotency key $\bullet$ replay $\bullet$ reconciliation $\bullet$
PCI DSS $\bullet$ SAQ-A
\end{keyterms}

\section{Three Ways to Take a Card, and Why One of Them Wins}

A shop that wants a card number has three architectures available, and the
choice determines how much of the payment industry's compliance regime applies
to it.

\begin{definitionbox}[Hosted page, Direct post, Server-to-server]
A \term{hosted page} sends the customer to the provider's own site to type
their card details, and receives them back afterwards with a result. The card
number never reaches your server.

\smallskip
A \term{direct post} keeps the customer on your page, but the form submits
straight to the provider. Your HTML is involved; your server is not.

\smallskip
\term{Server-to-server} means the card number arrives at your server and you
forward it. You now hold card data, and almost every obligation in this chapter
multiplies.

\smallskip
\textbf{For a shop of the size this course is about, the answer is always the
first or the second,} and the reason is in the next box.
\end{definitionbox}

\begin{conceptbox}[What PCI DSS actually asks of you]
The Payment Card Industry Data Security Standard governs anyone who stores,
processes or transmits card data. Its reporting burden is graded by how much of
that you do, through a Self-Assessment Questionnaire.

\begin{tight}
  \item \textbf{SAQ-A} is for a merchant who has \emph{entirely outsourced} card
        handling to a validated third party --- a hosted page. It is the
        shortest questionnaire, and most of it asks whether you are keeping
        your own site from interfering with the provider's.
  \item \textbf{SAQ-A-EP} applies to direct post: your page constructs the form,
        so your page can be tampered with, and you are asked more.
  \item \textbf{SAQ-D} is the full standard, and is what server-to-server costs
        you. It runs to hundreds of requirements.
\end{tight}

\smallskip
\textbf{The design decision and the compliance decision are the same
decision.} Choosing a hosted page is not a shortcut around PCI DSS; it is the
mechanism the standard offers for staying out of scope. A student who
understands only this from the chapter has understood the important part.

\smallskip
\textbf{What would change it:} storing the card number ``for the customer's
convenience'', logging the full request body, or accepting the number on your
page ``just to validate it''. Each of these moves a shop from SAQ-A to SAQ-D
with one line of code and no announcement.
\end{conceptbox}

\section{The Course Ships Its Own Gateway}

Measuring a real gateway needs a merchant account, which needs a registered
business. A sandbox needs credentials that cannot go in a handout and will have
expired before anyone reads it. So the stack in Appendix~A includes a provider
that does not exist.

\begin{alertbox}[What \texttt{code/gateway/provider.py} is, and is not]
It is not a toy card processor. It implements nothing about cards. It is a
gateway reduced to the \textbf{three properties that make integrating with one
difficult}, each configurable so that you can make it behave worse on purpose:

\begin{tight}
  \item \textbf{It is slow, and irregularly so.} A median around 300\,ms, with
        a minority of requests taking ten times that.
  \item \textbf{It sometimes does not answer at all.} About one request in
        fifty hangs, which is the case the rest of the chapter is about.
  \item \textbf{It tells you twice.} It honours an idempotency key if you send
        one, and charges twice if you do not.
\end{tight}

\smallskip
Those are the properties a real integration is actually built against.
Everything else a provider does --- 3-D Secure, tokenisation, settlement
reports --- is detail you will read in their documentation. These three are
what the \emph{architecture} has to survive, and they are the same for every
provider in every country.
\end{alertbox}

\section{Choosing a Timeout, Which Is Not a Performance Setting}

\begin{measured}[200 authorisations, waiting as long as it takes]
From \texttt{code/m14\_gateway.py --latency 200}, against the stack in
Appendix~A.

\smallskip
\begin{tabular}{@{}L{4.2cm}L{3.0cm}L{6.4cm}@{}}
Median & 321\,ms & what a single test request shows you \\
Mean & 700\,ms & \emph{more than twice the median} \\
90th percentile & 3{,}180\,ms & \\
95th percentile & 3{,}774\,ms & \\
99th percentile & 4{,}815\,ms & \\
Slowest answer & 4{,}998\,ms & \textbf{16$\times$ the median} \\
Never answered & 4 of 200 & 2\% \\
\end{tabular}

\smallskip
\textbf{The mean is more than twice the median, which is the whole
measurement.} A distribution where the mean sits far above the median is one
with a long tail, and an average response time of 700\,ms describes almost none
of the requests: most took a third of a second and a minority took five
seconds. Reporting ``average response time'' here would be worse than reporting
nothing, because it sounds like information.
\end{measured}

\begin{measured}[What each candidate timeout would have cost]
The same 200 requests, re-read against each timeout a developer might choose.

\rowcolors{2}{white}{lightbg}
\begin{longtable}{@{}L{2.4cm}L{2.4cm}L{2.2cm}L{2.4cm}L{4.2cm}@{}}
\toprule
\rowcolor{primary!15}
\textbf{Timeout} & \textbf{Answered} & \textbf{Cut off} & \textbf{Share} &
\textbf{What the shop is doing} \\
\midrule
\endhead
0.5\,s & 171 & 25 & 12.8\% & abandoning good payments \\
1.0\,s & 173 & 23 & 11.7\% & abandoning good payments \\
2.0\,s & 173 & 23 & 11.7\% & abandoning good payments \\
3.0\,s & 175 & 21 & 10.7\% & abandoning good payments \\
\textbf{5.0\,s} & \textbf{196} & \textbf{0} & \textbf{0.0\%} & waiting \\
8.0\,s & 196 & 0 & 0.0\% & waiting \\
\bottomrule
\end{longtable}

\smallskip
\textbf{The counter-intuitive result is the flat middle.} Raising the timeout
from 1\,second to 3\,seconds --- tripling how long the customer waits ---
rescues two payments in two hundred. One percent. Almost nothing.

\smallskip
The reason is that the distribution is not a bell curve but \emph{two groups}:
requests that answer in under half a second, and requests that take between two
and five seconds. There is hardly anything in between, so a timeout placed in
between catches hardly anything. \textbf{Either you wait for the slow group or
you do not, and every value between those two decisions buys you nothing while
costing the customer time.}

\smallskip
\textbf{A timeout chosen from the median rejects one payment in eight.} And
each rejection is not a failed payment --- it is a payment whose outcome the
shop does not know.
\end{measured}

\dsfig{%
\begin{tikzpicture}[
  st/.style={rounded corners=3pt, draw=#1, line width=1.1pt, fill=#1!8,
             text=#1!60!black, font=\bfseries\small, align=center,
             text width=3.1cm, minimum height=1.0cm},
  arr/.style={-{Stealth[length=2.2mm]}, primary!55, line width=0.9pt},
  note/.style={font=\scriptsize, text=neutral, align=left, text width=5.6cm,
               anchor=west}
]
  \node[st=secondary] (req) at (0,0) {Shop asks\\{\scriptsize\mdseries
       ``authorise this''}};
  \node[st=greenhl]   (ok)  at (5.0,1.5)  {Approved};
  \node[st=neutral]   (no)  at (5.0,0)    {Declined};
  \node[st=accent]    (un)  at (5.0,-1.5) {\textbf{No answer}};

  \foreach \t in {ok, no, un}{\draw[arr] (req) -- (\t);}

  \node[note, text=greenhl] at (6.9,1.5)
       {Capture, mark the order paid, dispatch.};
  \node[note, text=neutral] at (6.9,0)
       {Tell the customer, keep the order at status 0, let them retry.};
  \node[note, text=accent] at (6.9,-1.5)
       {\textbf{The money may or may not have left the customer's account,
        and the shop has no way to know.} Everything hard about payments is
        this one branch.};
  \node[font=\scriptsize\itshape, text=accent, align=center, text width=9.0cm]
       at (4.4,-3.1)
       {Two branches need code. The third needs a policy, a reconciliation
        process, and somebody's telephone number.};
\end{tikzpicture}}%
{The three outcomes of a payment request. Most integration tutorials cover the
first two.}{paymentoutcomes}

\begin{conceptbox}[What to do about the third branch]
A timed-out request has an unknown outcome. The shop must choose, in advance,
what it believes.

\begin{tightnum}
  \item \textbf{Never assume failure and never assume success.} Record it as
        \emph{unknown}, which is a third state and needs a place to live ---
        \chref{checkout}'s status 0 is where it goes.
  \item \textbf{Ask the provider.} Every gateway offers a status query. Retry
        it on a schedule, not in the customer's request.
  \item \textbf{Make the customer's retry safe.} If they press the button
        again, the same idempotency key must go with it, or you have charged
        them twice for being impatient.
  \item \textbf{Reconcile daily.} The provider's settlement report against your
        orders. This is the only mechanism that catches everything the other
        three miss, and it is the one small shops skip.
\end{tightnum}
\end{conceptbox}

\section{Idempotency, Demonstrated}

\begin{measured}[The same payment, sent twice]
From \texttt{code/m14\_gateway.py --idempotency}.

\smallskip
\textbf{Without a key:}

\smallskip
\begin{tabular}{@{}L{3.0cm}L{5.0cm}L{2.6cm}L{3.0cm}@{}}
first & transaction \texttt{f59594f7d23c4446} & approved & 125\,ms \\
second & transaction \texttt{8171b6e85f9b4766} & approved & 2{,}644\,ms \\
\end{tabular}

\smallskip
\textbf{Two transactions. The customer paid twice}, and both payments are
valid --- there is nothing for the provider to detect, because two requests to
charge the same amount are exactly what a customer buying two of something
looks like.

\smallskip
\textbf{With a key:}

\smallskip
\begin{tabular}{@{}L{3.0cm}L{5.0cm}L{2.6cm}L{3.0cm}@{}}
first & transaction \texttt{50f018abfd9243ec} & approved & 149\,ms \\
second & transaction \texttt{50f018abfd9243ec} & \textbf{replayed} &
\textbf{5\,ms} \\
\end{tabular}

\smallskip
The second request returned \emph{the first request's answer}, in 5\,ms,
without charging anything. The provider saw four requests and charged three
times.

\smallskip
\textbf{Where the key must come from.} Not from the request --- a retry is a
new request and would generate a new key, which defeats the entire mechanism.
It must come from \emph{the order}: generated once when the order row is
created, stored on it, and sent unchanged with every attempt to pay for that
order. This is the second reason \chref{checkout}'s order row exists before the
payment, and the better one.
\end{measured}

\begin{pitfall}
\begin{tight}
  \item \textbf{Choosing a timeout from the median response time.} It rejected
        one payment in eight here.
  \item \textbf{Reporting a gateway's ``average'' latency.} The mean was twice
        the median; the average describes no request that actually happened.
  \item \textbf{Treating a timeout as a failure} and letting the customer
        reorder.
  \item \textbf{Generating the idempotency key in the request} rather than
        storing it on the order.
  \item \textbf{Retrying inside the customer's request.} They are waiting; the
        retry belongs in a background job.
  \item \textbf{Logging the full request body} while integrating. If the body
        ever contains a card number, your log is now in PCI scope, and logs are
        backed up, shipped and searched by people who are not you.
  \item \textbf{Trusting a callback because it arrived.} Verify its signature;
        anyone can post to a public URL.
  \item \textbf{Skipping daily reconciliation} because the integration works.
        The integration working is what reconciliation is for.
\end{tight}
\end{pitfall}

\begin{lab}[Integrate against a gateway that misbehaves]
\begin{lstlisting}[language=bash]
# --- 1. Start the provider and look at it ---------------------------
cd code
docker compose up -d gateway
curl -s http://localhost:8092/health

# --- 2. Read it. It is one file -------------------------------------
# gateway/provider.py. Find the three properties it models, and the
# environment variable that controls each.

# --- 3. Measure before you decide -----------------------------------
python m14_gateway.py --latency 200

# --- 4. Predict your timeout, then check it -------------------------
# Write down the timeout you would have chosen from the median ALONE.
python m14_gateway.py --timeout

# --- 5. Demonstrate the double charge -------------------------------
python m14_gateway.py --idempotency

# --- 6. Make the provider worse, and re-decide ----------------------
# In docker-compose.yml set GATEWAY_TAIL_RATE to 0.35 and
# GATEWAY_HANG_RATE to 0.10, then:
docker compose up -d gateway
python m14_gateway.py --latency 100
# Does your timeout still hold? What changed, the median or the tail?

# --- 7. Put it back -------------------------------------------------
# Restore the defaults and restart the gateway.

# --- 8. Write the integration policy --------------------------------
# One page, four headings: timeout and why; what an unknown outcome
# means for the order; where the idempotency key comes from; how
# reconciliation happens and who does it.
\end{lstlisting}

\smallskip
\textbf{What to notice.} Step 6 is the point. Most people's timeout is chosen
once and never revisited, and a provider's tail is exactly the thing that
changes without telling you --- at the end of the month, on sale days, whenever
their own upstream is busy.
\end{lab}

\begin{checkpoint}
\begin{tightnum}
  \item A gateway's median response is 321\,ms and its 95th percentile is
        3{,}774\,ms. A colleague proposes a 1-second timeout, reasoning that it
        is ``three times the median, which is generous''. Using the measured
        table, say what that choice actually costs and what you would propose
        instead, with your reason.
  \item Explain why the measured timeout table barely improves between 1 second
        and 3 seconds, and what that shape tells you about the gateway that a
        single average could not.
  \item Your shop charges a customer, the request times out, and the customer
        presses Pay again. Describe exactly what happens with and without an
        idempotency key, say where the key must be stored, and explain why
        generating it when the button is pressed would not work.
\end{tightnum}
\end{checkpoint}

\begin{chaptersummary}
\begin{tight}
  \item A hosted page, a direct post and a server-to-server integration are the
        three architectures, and they map onto SAQ-A, SAQ-A-EP and SAQ-D. The
        design decision and the compliance decision are the same decision.
  \item A gateway is three hard properties: it is slow and irregular, it
        sometimes does not answer, and it tells you twice. The course ships a
        provider that does all three on purpose.
  \item Measured: median 321\,ms, mean 700\,ms, p95 3{,}774\,ms, slowest
        4{,}998\,ms --- sixteen times the median --- and 2\% that never
        answered at all. A mean twice the median is a tail, and an average is
        the wrong statistic for it.
  \item Measured: a 1-second timeout cuts off 11.7\% of payments and a
        3-second timeout cuts off 10.7\%. The distribution is two groups with a
        gap between them, so a timeout in the gap buys nothing.
  \item A timed-out payment is not a failed payment. It is an unknown outcome,
        it needs a state of its own, and it is resolved by querying the
        provider and by daily reconciliation --- not by guessing.
  \item Measured: without an idempotency key the same payment charged twice,
        with two valid transaction identifiers. With one, the second request
        returned the first answer in 5\,ms and charged nothing.
  \item The key belongs to the order, not to the request. This is the better
        reason why the order row exists before the payment.
\end{tight}
\end{chaptersummary}

\begin{reviewq}
  \item \textbf{(MCQ)} A shop using a hosted payment page completes which
        self-assessment questionnaire? (a)~SAQ-A; (b)~SAQ-A-EP; (c)~SAQ-D;
        (d)~none is required.
  \item \textbf{(MCQ)} In the measured distribution, the mean was 700\,ms and
        the median 321\,ms. This indicates: (a)~a measurement error;
        (b)~a long tail of slow requests; (c)~that half the requests failed;
        (d)~that the gateway is overloaded.
  \item \textbf{(Short)} Explain why a timed-out payment request has an unknown
        outcome rather than a failed one, and give the consequence of assuming
        either.
  \item \textbf{(Short)} Give two things that would move a shop from SAQ-A to a
        heavier questionnaire, each of which a developer might do without
        realising.
  \item \textbf{(Short)} Why must a callback's signature be verified, given
        that the callback comes from the provider?
  \item \textbf{(Applied)} You are integrating a provider whose published
        median is 250\,ms. Describe the measurement you would run before
        choosing a timeout, how many samples you would need and why, which
        statistics you would report, and what you would do differently if the
        distribution turned out to be a smooth bell curve instead of two
        groups.
  \item \textbf{(Applied)} Design the reconciliation process for a shop taking
        200 orders a day. State what two records you compare, how you match
        them, the three kinds of mismatch you expect to find, what each one
        means, and who in the business acts on it.
\end{reviewq}
